\documentclass{article}
\usepackage[T1]{fontenc}
\usepackage{lmodern}
\usepackage{graphicx}
\usepackage{amsmath,amssymb}
\usepackage[a4paper,margin=2cm]{geometry}
\usepackage[absolute,overlay]{textpos}
\setlength{\TPHorizModule}{1mm}
\setlength{\TPVertModule}{1mm}
\usepackage[indonesian]{babel}
\usepackage{hyperref}
\setlength{\emergencystretch}{2em}
% unit: Halaman 3

\begin{document}

% === Halaman 3 (python/native) ===

Teori Aljabar Abstrak

• Sebelum membahas ECC, perlu dipahami beberapa konsep aljabar 

abstrak yang mendasarinya.

• Konsep aljabar abstrak:

 1. Grup (group)

 2. Cincin (ring)

 3. Medan (field)

 3. Galois field


3

Rinaldi Munir/II4021 Kriptografi

\end{document}
